% GENERATED FILE. Edit canonical lessons, then run npm run generate:books.
% canonical-source-hash: fnv1a64:59cb604b1140ea96
% canonical-lessons: TA-C451-valakam, TA-C451-kilo, TA-C451-canta, TA-C451-ilappitu, TA-C451-cetamatai

\chapter{At the Shopping Complex}
\label{ch:ta-at-the-shopping-complex}
\hlchaptermodality{451}

\begin{chapteropening}
\textbf{By the end of this chapter:} \emph{I can talk about a shopping complex, kilos, a subscription and compensation, and goods getting damaged.}

Complete the last lesson of chapter 451: I can talk about a shopping complex, kilos, a subscription and compensation, and goods getting damaged.
\end{chapteropening}

\section[\texorpdfstring{\ta{வளாகம்} (vaḷākam)}{வளாகம் (vaḷākam)}]{\ta{வளாகம்} (vaḷākam) — a complex; a campus}

\label{lesson:TA-C451-valakam}



\begin{quote}
Before the new one: say the Tamil for the cost of living, then the Tamil for a banknote.
\end{quote}

\subsection*{You'll want to know: \ta{வளாகம்}}
\textbf{\ta{வளாகம்}} — \emph{vaḷākam} — \textquotedblleft{}a complex; a campus\textquotedblright{}.

\begin{quote}
\ta{புதிய} \ta{வணிக} \ta{வளாகத்தில்} \ta{நிறைய} \ta{கடைகள்} \ta{இருக்கின்றன}.
\end{quote}

\begin{quote}
\emph{putiya vaṇika vaḷākattil niṟaiya kaṭaikaḷ irukkiṉṟaṉa.}
\end{quote}

\begin{quote}
There are many shops in the new shopping complex.
\end{quote}

\ta{வணிக} \ta{வளாகம்} is a shopping complex, and \ta{பல்கலைக்கழக} \ta{வளாகம்} a campus.

\begin{cousinweb}
\textbf{\ta{வளாகம்}} (\emph{vaḷākam}) is native Dravidian (DEDR 5313), \textquotedblleft{}enclosing, surrounding\textquotedblright{}, in the article of \ta{வளை}, \textquotedblleft{}to surround\textquotedblright{}.
\end{cousinweb}

\begin{grammarlens}[title={What you've built}]
One more word for getting around.
\end{grammarlens}

\subsection*{Your turn}
\begin{itemize}
\raggedright
  \item \emph{Say it:} \emph{vaḷākam}
  \item \emph{Say it:} \emph{vaḷākam}, once more
  \item \emph{Say it:} \emph{putiya vaṇika vaḷākattil niṟaiya kaṭaikaḷ irukkiṉṟaṉa.} — There are many shops in the new shopping complex.
  \item \emph{Recall:} say the Tamil for the cost of living, then the Tamil for a banknote, then say \emph{vaḷākam} again
\end{itemize}

\begin{tcolorbox}[breakable,colback=teal!4,colframe=teal!35!black,title={Before you move on}]
What does \ta{வளாகம்} mean? (\textquotedblleft{}A complex; a campus\textquotedblright{}.) Say it once more.
\end{tcolorbox}
\section[\texorpdfstring{\ta{கிலோ} (kilō)}{கிலோ (kilō)}]{\ta{கிலோ} (kilō) — a kilo}

\label{lesson:TA-C451-kilo}



\begin{quote}
Before the new one: say the Tamil for a banknote, then the Tamil for a complex.
\end{quote}

\subsection*{You'll want to know: \ta{கிலோ}}
\textbf{\ta{கிலோ}} — \emph{kilō} — \textquotedblleft{}a kilo\textquotedblright{}.

\begin{quote}
\ta{ஒரு} \ta{கிலோ} \ta{தக்காளி} \ta{எவ்வளவு}?
\end{quote}

\begin{quote}
\emph{oru kilō takkāḷi evvaḷavu?}
\end{quote}

\begin{quote}
How much is a kilo of tomatoes?
\end{quote}

Numbers come before it: \ta{இரண்டு} \ta{கிலோ} \ta{அரிசி}.

\begin{cousinweb}
\textbf{\ta{கிலோ}} (\emph{kilō}) is a loan from English \emph{kilo}.
\end{cousinweb}

\begin{grammarlens}[title={What you've built}]
One more word for getting around.
\end{grammarlens}

\subsection*{Your turn}
\begin{itemize}
\raggedright
  \item \emph{Say it:} \emph{kilō}
  \item \emph{Say it:} \emph{kilō}, once more
  \item \emph{Say it:} \emph{oru kilō takkāḷi evvaḷavu?} — How much is a kilo of tomatoes?
  \item \emph{Recall:} say the Tamil for a banknote, then the Tamil for a complex, then say \emph{kilō} again
\end{itemize}

\begin{tcolorbox}[breakable,colback=teal!4,colframe=teal!35!black,title={Before you move on}]
What does \ta{கிலோ} mean? (\textquotedblleft{}A kilo\textquotedblright{}.) Say it once more.
\end{tcolorbox}
\section[\texorpdfstring{\ta{சந்தா} (cantā)}{சந்தா (cantā)}]{\ta{சந்தா} (cantā) — a subscription; a membership fee}

\label{lesson:TA-C451-canta}



\begin{quote}
Before the new one: say the Tamil for a complex, then the Tamil for a kilo.
\end{quote}

\subsection*{You'll want to know: \ta{சந்தா}}
\textbf{\ta{சந்தா}} — \emph{cantā} — \textquotedblleft{}a subscription; a membership fee\textquotedblright{}.

\begin{quote}
\ta{நாளிதழுக்கு} \ta{மாதச்} \ta{சந்தா} \ta{கட்டினேன்}.
\end{quote}

\begin{quote}
\emph{nāḷitaḻukku mātac cantā kaṭṭiṉēṉ.}
\end{quote}

\begin{quote}
I paid the monthly subscription for the newspaper.
\end{quote}

\ta{சந்தா} \ta{கட்டு} is to pay a subscription.

\begin{cousinweb}
\textbf{\ta{சந்தா}} (\emph{cantā}) is usually traced to Hindustani \emph{candā}, \textquotedblleft{}a subscription, a contribution\textquotedblright{}; the path into Tamil is not settled.
\end{cousinweb}

\begin{grammarlens}[title={What you've built}]
One more word for getting around.
\end{grammarlens}

\subsection*{Your turn}
\begin{itemize}
\raggedright
  \item \emph{Say it:} \emph{cantā}
  \item \emph{Say it:} \emph{cantā}, once more
  \item \emph{Say it:} \emph{nāḷitaḻukku mātac cantā kaṭṭiṉēṉ.} — I paid the monthly subscription for the newspaper.
  \item \emph{Recall:} say the Tamil for a complex, then the Tamil for a kilo, then say \emph{cantā} again
\end{itemize}

\begin{tcolorbox}[breakable,colback=teal!4,colframe=teal!35!black,title={Before you move on}]
What does \ta{சந்தா} mean? (\textquotedblleft{}A subscription; a membership fee\textquotedblright{}.) Say it once more.
\end{tcolorbox}
\section[\texorpdfstring{\ta{இழப்பீடு} (iḻappīṭu)}{இழப்பீடு (iḻappīṭu)}]{\ta{இழப்பீடு} (iḻappīṭu) — compensation}

\label{lesson:TA-C451-ilappitu}



\begin{quote}
Before the new one: say the Tamil for a kilo, then the Tamil for a subscription.
\end{quote}

\subsection*{You'll want to know: \ta{இழப்பீடு}}
\textbf{\ta{இழப்பீடு}} — \emph{iḻappīṭu} — \textquotedblleft{}compensation\textquotedblright{}.

\begin{quote}
\ta{விமானம்} \ta{தாமதமானதால்} \ta{பயணிகளுக்கு} \ta{இழப்பீடு} \ta{கொடுத்தார்கள்}.
\end{quote}

\begin{quote}
\emph{vimāṉam tāmatamāṉatāl payaṇikaḷukku iḻappīṭu koṭuttārkaḷ.}
\end{quote}

\begin{quote}
Because the flight was late, they gave the passengers compensation.
\end{quote}

\ta{இழப்பு}, a loss, + īṭu, making good.

\begin{cousinweb}
\textbf{\ta{இழப்பீடு}} (\emph{iḻappīṭu}) is built from \ta{இழப்பு}, \textquotedblleft{}a loss\textquotedblright{} (DEDR 501, the family of \ta{இழ}, to lose), and \emph{īṭu}, \textquotedblleft{}putting in, making good\textquotedblright{} (DEDR 442).
\end{cousinweb}

\begin{grammarlens}[title={What you've built}]
One more word for getting around.
\end{grammarlens}

\subsection*{Your turn}
\begin{itemize}
\raggedright
  \item \emph{Say it:} \emph{iḻappīṭu}
  \item \emph{Say it:} \emph{iḻappīṭu}, once more
  \item \emph{Say it:} \emph{vimāṉam tāmatamāṉatāl payaṇikaḷukku iḻappīṭu koṭuttārkaḷ.} — Because the flight was late, they gave the passengers compensation.
  \item \emph{Recall:} say the Tamil for a kilo, then the Tamil for a subscription, then say \emph{iḻappīṭu} again
\end{itemize}

\begin{tcolorbox}[breakable,colback=teal!4,colframe=teal!35!black,title={Before you move on}]
What does \ta{இழப்பீடு} mean? (\textquotedblleft{}Compensation\textquotedblright{}.) Say it once more.
\end{tcolorbox}
\section[\texorpdfstring{\ta{சேதமடை} (cētamaṭai)}{சேதமடை (cētamaṭai)}]{\ta{சேதமடை} (cētamaṭai) — to get damaged}

\label{lesson:TA-C451-cetamatai}



\begin{quote}
Before the new one: say the Tamil for a subscription, then the Tamil for compensation.
\end{quote}

\subsection*{You'll want to know: \ta{சேதமடை}}
\textbf{\ta{சேதமடை}} — \emph{cētamaṭai} — \textquotedblleft{}to get damaged\textquotedblright{}.

\begin{quote}
\ta{வெள்ளத்தில்} \ta{பல} \ta{வீடுகள்} \ta{சேதமடைந்தன}.
\end{quote}

\begin{quote}
\emph{veḷḷattil pala vīṭukaḷ cētamaṭaintaṉa.}
\end{quote}

\begin{quote}
Many houses were damaged in the flood.
\end{quote}

Past (it …): \ta{சேதமடைந்தது}, \emph{cētamaṭaintatu}.

\ta{சேதம்}, damage, + \ta{அடை}, to get.

\begin{cousinweb}
\textbf{\ta{சேதமடை}} (\emph{cētamaṭai}) is built from \ta{சேதம்}, \textquotedblleft{}damage\textquotedblright{}, borrowed from Sanskrit \emph{cheda}, \textquotedblleft{}cutting, loss\textquotedblright{} (CDIAL 5064), and \ta{அடை}, \textquotedblleft{}to reach, to get\textquotedblright{} (DEDR 79).
\end{cousinweb}

\begin{grammarlens}[title={What you've built}]
One more word for getting around.
\end{grammarlens}

\subsection*{Your turn}
\begin{itemize}
\raggedright
  \item \emph{Say it:} \emph{cētamaṭai}
  \item \emph{Say it:} \emph{cētamaṭai}, once more
  \item \emph{Say it:} \emph{veḷḷattil pala vīṭukaḷ cētamaṭaintaṉa.} — Many houses were damaged in the flood.
  \item \emph{Recall:} say the Tamil for a subscription, then the Tamil for compensation, then say \emph{cētamaṭai} again
\end{itemize}

\begin{tcolorbox}[breakable,colback=teal!4,colframe=teal!35!black,title={Before you move on}]
What does \ta{சேதமடை} mean? (\textquotedblleft{}To get damaged\textquotedblright{}.) Say it once more.
\end{tcolorbox}
